\section{De «descargar internet» a «descargar a los humanos / el mundo»}

Una de las metáforas que más se difundieron de la entrevista es la del «OpenAI inverso». La ruta directa es: descargar datos de internet, entrenar un transformer / GPT, obtener inteligencia lingüística y después llevarla al mercado. Esta ruta funciona porque internet ya contiene una gran cantidad de conocimiento escrito por personas. Los modelos del mundo no cuentan con el mismo shortcut: los datos del mundo real, los escenarios empresariales, los procesos físicos, las tareas robóticas y la retroalimentación de sensores no pueden simplemente descargarse de páginas web.

\begin{figure}[H]
\centering
\includegraphics[width=0.96\textwidth]{download-internet-to-world.png}
\caption{De «descargar internet» a «descargar a los humanos / el mundo»: los modelos del mundo necesitan entornos reales y un ciclo cerrado con socios.}
\end{figure}

\begin{knowledgebox}{Lectura de la figura: lo esencial de la ruta inversa es el ciclo cerrado de colaboración}
A la izquierda está la ruta directa, de internet a los modelos de lenguaje; a la derecha, la ruta inversa, en la que el mundo, los socios, las tareas y los datos retroalimentan en conjunto al modelo. Las condiciones del ciclo cerrado en la parte inferior de la figura indican que un modelo del mundo necesita datos de socios, un modelo inicial, tareas reales, retroalimentación de evaluación y control de seguridad. Sin la participación del mundo, el modelo del mundo se queda en un concepto de artículo académico o en una demo generativa.
\end{knowledgebox}

\subsection{World model needs the world}

Xie Saining (谢赛宁) dice: world model needs the world. Esta frase es el núcleo de la estrategia organizacional de AMI Labs. Significa que el modelo del mundo no es una tarea que una sola empresa pueda completar con datos cerrados, sino que requiere la colaboración de granjas, hospitales, fábricas, empresas de robótica, redes de sensores, entornos de simulación y escenarios de negocio reales.

En este marco, el modelo aporta primero una capacidad inicial y entra en escenarios reales para generar valor; los escenarios reales producen retroalimentación, errores y datos nuevos; y esos datos retroalimentan a su vez al modelo. Este ciclo cerrado se parece a un motor de datos, pero su objeto ya no es solo el texto de páginas web, sino el mundo físico y los sistemas de acción.

\begin{importantbox}{La esencia del OpenAI inverso}
El OpenAI directo descarga internet; el OpenAI inverso tiene que construir los datos junto con el mundo. El primero depende de corpus de texto ya existentes; el segundo depende de tareas reales, redes de socios, sistemas de evaluación y retroalimentación continua.
\end{importantbox}

\subsection{Por qué esto no es simplemente comprar datos}

Decir que world model needs the world puede hacer pensar que la solución consiste solo en adquirir más datos físicos. Esta sección desmonta ese malentendido: lo que de verdad necesita un modelo del mundo es un mecanismo de generación de datos, no un inventario de datos de una sola vez.

Si lo que necesita el modelo del mundo son las consecuencias de cada action, los estados físicos, la recuperación ante fallas, los flujos de sensores y los procesos de cada dominio, entonces «comprar un lote de datos» está muy lejos de ser suficiente. Los datos deben tener tareas, objetivos, retroalimentación, evaluación y actualización continua. Más importante aún, muchos datos solo aparecen después de que el modelo entra en el escenario: el modelo se equivoca, una persona lo corrige, el sistema se recupera, el usuario cambia el flujo de trabajo; nada de esto puede prepararse de antemano en un corpus estático.

\begin{knowledgebox}{De data factory a environment factory}
Este episodio puede leerse junto con el panorama de datos de EP134: lo que necesita un modelo del mundo no son datos adquiridos en una sola compra, sino entornos capaces de producir de forma continua tareas, fallas, retroalimentación y evaluaciones. La fábrica de datos del futuro se parecerá más a una environment factory.
\end{knowledgebox}

\subsection{Resumen de la sección}

«Descargar internet» resolvió la primera etapa de la inteligencia lingüística; «descargar a los humanos / el mundo» implica pasar del texto estático a los entornos dinámicos. La competencia en modelos del mundo no es solo una competencia de parámetros del modelo, sino también de redes de colaboración, sistemas de evaluación y ciclos cerrados de retroalimentación real.
